\uuid{albb}
\titre{Limite}
\auteur{bodin ; hanani ; mzari}
\organisation{exo7}
\licence{CC BY-NC-SA 4.0}
\datecreate{2020-01-20}
\niveau{L1}
\module{Analyse}
\chapitre{Suite}
\sousChapitre{Autre}
\difficulte{2}
\theme{Suites, limite}
\exo7qcmid{lille-1.suites.013}

\contenu{
\begin{question}
\qclassification{121.00}
\qtags{motcle=limite}
Soit $\displaystyle u_n=\frac{2n^2+1}{n^2+1}$. Quelles sont les bonnes réponses ?
\begin{answers}
    \bad{$\forall \varepsilon >0,\; \forall n\in \Nn,\; |u_n-2|<\varepsilon$}
    \good{$\exists \varepsilon >0,\; \forall n\in \Nn,\; |u_n-2|<\varepsilon$}
    \good{$\forall n\in \Nn,\; n>10\Rightarrow |u_n-2|<10^{-2}$}
    \good{$\forall \varepsilon >0,\; \exists n_0\in \Nn,\; \forall n\in \Nn,\; n>n_0\Rightarrow |u_n-2|<\varepsilon$}
\end{answers}
\begin{explanations}
D'abord, $\displaystyle u_n-2=\frac{-1}{n^2+1}$. D'une part, $\displaystyle |u_n-2|\leq 1$ pour tout $n\in \Nn$. D'autre part, si $n>10$ alors $\displaystyle n^2+1>10^2+1>10^2$. C'est-à-dire
$$|u_n-2|=\frac{1}{n^2+1}<10^{-2}.$$
De même, pour tout $\varepsilon >0$, si $n>n_0=E\left(\sqrt{\varepsilon ^{-1}}\right)+1$ alors $\displaystyle |u_n-2|=\frac{1}{n^2+1}<\varepsilon$.
\end{explanations}
\end{question}
}
